\begin{table}[tbp]
\centering
\footnotesize
\setlength{\tabcolsep}{3.2pt}
\renewcommand{\arraystretch}{1.12}
\caption{Properties enforced by the typed support relation.}
\label{tab:s-support-properties}
\begin{tabularx}{\textwidth}{@{}l >{\raggedright\arraybackslash}p{0.25\textwidth} Y@{}}
\toprule
\textbf{ID} & \textbf{Property} & \textbf{Operational consequence} \\
\midrule
P1 & Scope monotonicity & Widening a claim adds requirements; it cannot erase requirements already named. \\
P2 & Evidence restriction & Removing the only admissible witness for a required element breaks support. \\
P3 & Rank invariance & Strictly increasing score calibration preserves AUROC/AUPRC ordering but may change a fixed-threshold decision. \\
P4 & Protocol non-substitutability & A result under one claim-relevant visibility or selection coordinate cannot fill another without an invariance argument. \\
P5 & Resource separation & A resource-blocked cell is measurable for feasibility but unordered by predictive performance. \\
\bottomrule
\end{tabularx}
\vspace{2pt}
\begin{minipage}{\textwidth}\footnotesize\emph{Scope and reading.} These are specification and ordering properties, not claims that the ontology is complete or that the validator establishes scientific truth.\end{minipage}
\end{table}
